% Ad Atticum 7.20 (ad-atticum-07-20)
% Source: https://raw.githubusercontent.com/PerseusDL/canonical-latinLit/master/data/phi0474/phi057/phi0474.phi057.perseus-lat1.xml
% Fetched by scripts/fetch_latin.py

% Dateline (Perseus): Scr. Capuae Non. Febr. a. 705 (49) .

\ciceroLetterOpener{CICERO ATTICO salutem}

\ciceroSection{1} breviloquentem iam me tempus ipsum facit. pacem enim desperavi, bellum nostri nullum administrant. cave enim putes quicquam esse minoris his consulibus; quorum ego spe audiendi aliquid et cognoscendi nostri apparatus maximo imbri Capuam veni pridie Nonas, ut eram iussus. illi autem nondum venerant sed erant venturi inanes, imparati. Gnaeus autem Luceriae dicebatur esse et adire cohortis legionum Appianarum non firmissimarum. at illum ruere nuntiant et iam iamque adesse, non ut manum conserat (quicum enim?) sed ut fugam intercludat.

\ciceroSection{2} ego autem in Italia καὶ συναποθανεῖν —nec te id consulo; sin extra, quid ago? ad manendum hiems, lictores, improvidi et neglegentes duces, ad fugam hortatur amicitia Gnaei, causa bonorum, turpitudo coniungendi cum tyranno; qui quidem incertum est Phalarimne an Pisistratum sit imitaturus. haec velim explices et me iuves consilio; etsi te ipsum istic iam calere puto, sed tamen quantum poteris. ego si quid hic hodie novi cognoro, scies; iam enim aderunt consules ad suas Nonas. tuas cotidie litteras exspectabo; ad has autem cum poteris rescribes. mulieres et Cicerones in Formiano reliqui.
